% Vietnamese translation for OI Public Library
% Source-File: IOI中国国家候选队论文/2003/饶向荣/饶向荣.ppt
\documentclass[11pt]{article}
\usepackage{oipl-vn}

\title{Bản trình chiếu: DNA của virus}
\author{Rao Xiangrong}
\date{2003}

\begin{document}
\maketitle

\origtitle{Virus DNA: analyzing the solution process of a string matching problem}
\sourcepath{IOI China candidate team papers / 2003 / Rao Xiangrong / Rao Xiangrong.ppt}

\section{Mô tả bài toán}

Bài trình bày xét một bài khớp chuỗi. Virus có chuỗi DNA dạng vòng, trong khi
DNA của sinh vật là chuỗi thẳng. Khả năng sinh vật bị virus xâm nhiễm phụ
thuộc vào độ dài phần chung dài nhất giữa DNA sinh vật và DNA virus. Vì DNA
virus là vòng, nó có thể khớp lặp tuần hoàn.

Nếu phần chung dài nhất ngắn hơn độ dài DNA virus, xác suất xâm nhiễm được
xem là \(0\). Ngược lại, slide cho công thức:
\[
  \text{xác suất xâm nhiễm}
  =
  \frac{\text{độ dài phần chung dài nhất}}{\text{độ dài DNA sinh vật}}.
\]

Gọi \(A\) là chuỗi vòng của virus, độ dài \(N\), và \(B\) là chuỗi thẳng của
sinh vật, độ dài \(M\). Bài toán cần tìm độ dài chuỗi con chung dài nhất giữa
chuỗi vòng \(A\), cho phép lặp tuần hoàn, và chuỗi thẳng \(B\). Ràng buộc
trong slide là
\[
  N\le 1000,\qquad M\le 10^5.
\]
Ví dụ: nếu \(A=\texttt{abc}\) và \(B=\texttt{abbcabcabb}\), phần chung dài
nhất có độ dài \(7\), là \(\texttt{bcabcab}\).

\section{Quy hoạch động trực tiếp}

Cách nghĩ đầu tiên là dùng quy hoạch động cho chuỗi con chung dài nhất. Gọi
\[
  f[i,j]
\]
là độ dài chuỗi con chung dài nhất kết thúc tại vị trí \(i\) của \(B\) và vị
trí \(j\) của \(A\). Khi \(B[i]=A[j]\), trạng thái kéo dài từ cặp trước đó;
khi hai ký tự khác nhau, giá trị tương ứng bằng \(0\). Vì \(A\) là chuỗi
vòng, chỉ số trên \(A\) được hiểu theo modulo \(N\).

Không cần tính mọi \(f[i,j]\): chỉ khi \(B[i]=A[j]\) mới có khả năng có giá
trị khác \(0\). Vì bảng chữ cái trong slide là các chữ cái
\(\texttt{a}\ldots\texttt{z}\) và \(\texttt{A}\ldots\texttt{Z}\), có thể lưu
danh sách vị trí xuất hiện của từng ký tự trong \(A\) để chỉ duyệt các \(j\)
phù hợp. Không gian có thể giảm còn \(O(N)\).

Nhược điểm là thời gian vẫn là \(O(MN)\) trong trường hợp xấu nhất. Với
\(N\le1000\) và \(M\le10^5\), slide đánh giá độ phức tạp này là quá cao.

\section{Điều kiện then chốt}

Điều kiện chưa được DP trực tiếp khai thác là chỉ khi độ dài phần chung đạt
ít nhất \(N\) thì mới cần tính xác suất khác \(0\). Nếu độ dài tối đa nhỏ hơn
\(N\), có thể trả lời \(0\) ngay.

Do vị trí bắt đầu trên chuỗi vòng không cố định, slide cắt vòng \(A\) tại một
vị trí bất kỳ thành chuỗi thẳng \(C\). Nếu phần chung có độ dài ít nhất
\(N\), phần chung đó phải chứa đủ một vòng của \(C\). Vì vậy bài toán được
đưa về các phép khớp từ \(B\) với \(C\), nhưng cho phép sau khi khớp hết
\(C\) thì tiếp tục lặp lại \(C\).

Định nghĩa:
\begin{itemize}
  \item \(R[i]\): độ dài lớn nhất khi bắt đầu ở \(B[i]\) và \(C[1]\), khớp
  về phía sau theo chu kỳ của \(C\);
  \item \(L[i-1]\): độ dài lớn nhất khi bắt đầu ở \(B[i-1]\) và \(C[N]\),
  khớp ngược về phía trước.
\end{itemize}
Các mảng này cho phép ghép phần khớp qua điểm cắt của chuỗi vòng.

\section{Từ \texorpdfstring{\(R[i]\)}{R[i]} đến \texorpdfstring{\(Q[i]\)}{Q[i]}}

Cách tính thô \(R[i]\) là thử từng vị trí \(i\) trong \(B\), rồi so sánh
từng ký tự với \(C\) cho đến khi không khớp. Độ dài khớp có thể đến \(M\),
nên cách này có thể thành \(O(M^2)\).

Slide quan sát rằng nếu từ vị trí \(i\) đã khớp được ít nhất \(N\) ký tự, thì
không cần tiếp tục khớp một cách mù quáng:
\[
  R[i]=R[i+N]+N.
\]
Do đó chỉ cần biết độ dài khớp không tuần hoàn giữa hậu tố của \(B\) và
chuỗi \(C\). Định nghĩa
\[
  Q[i]
\]
là độ dài lớn nhất khi bắt đầu ở \(B[i]\) và \(C[1]\), nhưng chỉ khớp
\(C\) như một chuỗi thẳng không lặp. Tính được \(Q[i]\) thì có thể suy ra
\(R[i]\) bằng công thức lặp theo chu kỳ.

Kế hoạch dùng trực tiếp KMP để lấy mọi \(Q[i]\) không thuận tiện, vì quá
trình KMP nhảy vị trí và không cho ngay độ dài khớp bắt đầu tại mọi \(i\).
Tuy vậy, ý tưởng chính của KMP vẫn hữu ích: tận dụng kết quả so sánh trước
đó để không so sánh lại các ký tự của \(B\) đã biết.

\section{Tính \texorpdfstring{\(Q[i]\)}{Q[i]} bằng vùng đã khớp}

Giả sử trong quá trình tính \(Q[1],Q[2],\ldots,Q[i-1]\), vị trí xa nhất của
\(B\) đã được khớp tới là \(k\), và lần khớp đạt tới \(k\) bắt đầu tại vị trí
\(j\). Khi xét \(Q[i]\), ta muốn bắt đầu so sánh mới từ \(k+1\) nếu chắc chắn
phần từ \(i\) đến \(k\) đã khớp.

Đoạn \(B[j..k]\) giống với đoạn \(C[1..k-j+1]\). Vì vị trí \(i\) trong
\(B\) tương ứng với vị trí \(i-j+1\) trong \(C\), ta cần biết từ
\(C[i-j+1]\) và \(C[1]\) có thể khớp được bao lâu. Gọi độ dài đó là \(L\).
Khi đó:
\[
  L<k-i+1 \Rightarrow Q[i]=L,
\]
còn nếu
\[
  L\ge k-i+1,
\]
thì \(Q[i]+i\ge k\), và ta chỉ cần tiếp tục so sánh từ \(k+1\) để tính
\(Q[i]\).

Như vậy điểm then chốt là tính nhanh các giá trị \(L\) chỉ phụ thuộc vào
chuỗi \(C\).

\section{Mảng \texorpdfstring{\(T\)}{T} và độ phức tạp}

Đặt
\[
  T[x]\qquad (1\le x\le N)
\]
là độ dài lớn nhất khi bắt đầu ở \(C[x]\) và \(C[1]\). Trong bước tính
\(Q[i]\), giá trị cần dùng chính là
\[
  T[i-j+1].
\]

Việc tính \(T\) có bản chất giống tính \(Q\): đều cần biết từ mỗi vị trí của
một chuỗi có thể khớp với \(C\) được bao lâu. Vì vậy \(T[x]\) có thể được
tính dựa trên các giá trị \(T[1],\ldots,T[x-1]\) và một lượng so sánh bổ sung
thích hợp. Slide ghi tổng thời gian của toàn bộ thuật toán là
\[
  O(M+N),
\]
gồm thời gian khớp trên \(B\) và thời gian tính \(T\).

\section{Tổng kết}

Quá trình giải được slide tóm tắt như sau:
\begin{enumerate}
  \item thử quy hoạch động trực tiếp, nhưng thời gian không đạt;
  \item phân tích đề và nhận ra điều kiện quan trọng là độ dài cần xét phải
  không nhỏ hơn \(N\);
  \item cắt chuỗi vòng thành chuỗi thẳng \(C\), rồi tách lời giải qua các
  mảng \(L[i]\) và \(R[i]\);
  \item rút gọn việc tính \(R[i]\) về \(Q[i]\);
  \item mượn tư tưởng của KMP để thiết kế mảng \(T\), giảm thời gian khớp
  xuống tuyến tính.
\end{enumerate}

Slide nhấn mạnh rằng lời giải này giống KMP về tư tưởng tránh so sánh lặp,
nhưng không phải là áp dụng máy móc thuật toán KMP. Bài học chính là cần phân
tích đầy đủ điều kiện của đề, nắm chắc kiến thức nền về khớp chuỗi, và linh
hoạt biến đổi mô hình trước khi cố tối ưu một DP trực tiếp.

\end{document}
